\documentclass[11pt,a4paper]{article}

\usepackage[margin=1in]{geometry}
\usepackage[T1]{fontenc}
\usepackage[utf8]{inputenc}
\usepackage{lmodern}
\usepackage{amsmath,amssymb,amsthm,mathtools}
\usepackage{bm}
\usepackage{booktabs}
\usepackage{array}
\usepackage{enumitem}
\usepackage{hyperref}
\usepackage{xcolor}

\hypersetup{
  colorlinks=true,
  linkcolor=blue!50!black,
  urlcolor=blue!50!black,
  citecolor=blue!50!black
}

\title{Mathematical Ideas Tested in the \texttt{Diss\_eng} Repository}
\author{Repository analysis document}
\date{\today}

\begin{document}
\maketitle

\begin{abstract}
This repository is not a single monolithic codebase; it is a small research workspace built around one central theme: how well approximate quantum state preparation can be improved, analyzed, or structured by mathematically designed dissipative dynamics and by target-aware state-loading methods. The main mathematical content appears in two experiment families. The first, \path{diss_lindblad_tail/}, studies density-matrix evolution under Lindblad master equations after a truncated state-preparation circuit. The second, \path{experiments/hydrogen_qalchemy/}, constructs hydrogenic wavefunctions on a discretized spatial grid, maps them to quantum statevectors, prepares them through Q-Alchemy-generated circuits, and then studies both ideal and hardware-constrained dissipative tails. A third component, \path{q-alchemy-sdk-py/}, is mostly infrastructural, but it reflects the mathematics of approximate quantum state preparation, tensor-inspired initialization methods, and fidelity-constrained circuit synthesis. This document explains the mathematical structures actually exercised in the repository and the kinds of hypotheses the code is testing.
\end{abstract}

\section{Repository-Level Research Question}
At a high level, the repository tests variants of the following question:
\begin{quote}
Given a target pure quantum state $\lvert \psi_\star\rangle$, and given either an approximate state-preparation circuit or a restricted hardware model, can one design dissipative dynamics that provably or empirically increase fidelity to $\lvert \psi_\star\rangle$ while preserving physicality of the quantum state?
\end{quote}

That question is explored in two progressively richer settings:
\begin{enumerate}[label=\arabic*.]
  \item An \emph{exact} open-system simulation setting, where the dissipative dynamics are represented by a full Liouvillian superoperator and evolved deterministically in density-matrix form.
  \item A \emph{large-scale and hardware-constrained} setting, where the target state is a discretized hydrogenic wavefunction and the dissipative tail is approximated by local Kraus channels, edge channels, and ancilla-reset gadgets.
\end{enumerate}

The relevant modules are concentrated in:
\begin{itemize}
  \item \path{diss_lindblad_tail/diss_lindblad/}
  \item \path{experiments/hydrogen_qalchemy/}
  \item \path{experiments/hydrogen_qalchemy/tail/}
\end{itemize}

\section{Core Mathematical Objects}
Across the repository, the same mathematical objects reappear.

\subsection{Pure states and density matrices}
The basic target is a normalized statevector
\[
\lvert \psi \rangle \in \mathbb{C}^{2^n},
\qquad
\langle \psi \mid \psi \rangle = 1.
\]
Pure states are embedded into density-matrix language by
\[
\rho = \lvert \psi \rangle \langle \psi \rvert.
\]
This appears explicitly in \path{diss_lindblad_tail/diss_lindblad/density_matrix.py} and \path{experiments/hydrogen_qalchemy/tail/metrics.py}.

The repository repeatedly tests whether an evolution rule preserves the defining conditions of a physical density matrix:
\begin{align*}
\rho &= \rho^\dagger, \\
\operatorname{Tr}(\rho) &= 1, \\
\rho &\succeq 0.
\end{align*}
These are not only assumed; they are checked in the test suites.

\subsection{Target fidelity}
The main scalar performance metric is fidelity to a pure target:
\[
F(\rho,\psi_\star) = \langle \psi_\star \rvert \rho \lvert \psi_\star \rangle.
\]
This is the principal observable in both experiment families. The repository studies whether fidelity:
\begin{itemize}
  \item increases after adding a dissipative tail,
  \item follows a known analytic law,
  \item remains monotone under an ideal cooling channel,
  \item can be approximated by constrained channels fitted to a teacher model.
\end{itemize}

\subsection{Quantum channels, CPTP maps, and Lindblad generators}
The code alternates between two equivalent open-system descriptions:
\begin{enumerate}[label=\alph*.]
  \item continuous-time Lindblad evolution,
  \item discrete completely positive trace-preserving (CPTP) maps in Kraus form.
\end{enumerate}

The Lindblad master equation implemented in \path{diss_lindblad_tail/diss_lindblad/lindblad.py} is
\[
\frac{d\rho}{dt}
= -i[H,\rho] + \sum_k \gamma_k \mathcal{D}[L_k](\rho),
\]
where
\[
\mathcal{D}[L](\rho)
= L\rho L^\dagger - \frac{1}{2}\{L^\dagger L,\rho\}.
\]
This is then vectorized into Liouville space and exponentiated as a matrix evolution problem.

In the constrained-tail code, each short time step is instead implemented as a composition of local Kraus maps. The mathematical question there is whether an expressive enough family of physically realizable CPTP maps can mimic a desired global dissipative flow.

\subsection{State-preparation unitaries}
Another recurring idea is the extension of a target state to a full unitary:
\[
U \lvert 0\cdots 0\rangle = \lvert \psi_\star \rangle.
\]
In \path{diss_lindblad_tail/diss_lindblad/circuits.py}, the first column of a unitary is chosen to be $\psi_\star$, and the remaining columns are completed by an orthonormal basis from a null-space computation. The exact circuit is then decomposed into elementary gates and deliberately truncated. This means the repository is testing not only dissipative dynamics, but the interaction between:
\begin{itemize}
  \item exact state preparation,
  \item shallow or incomplete approximation,
  \item post-circuit dissipative correction.
\end{itemize}

\section{Mathematical Ideas in \texttt{diss\_lindblad\_tail/}}

\subsection{Approximate state preparation followed by a dissipative tail}
The package \path{diss_lindblad_tail/} implements a clean mathematical pipeline:
\begin{enumerate}[label=\arabic*.]
  \item choose a target state $\lvert \psi_\star\rangle$,
  \item build an exact preparation circuit,
  \item truncate the circuit to a gate fraction,
  \item convert the resulting approximate output into a density matrix $\rho_0$,
  \item evolve $\rho_0$ under a chosen dissipative model,
  \item record fidelity, trace, and purity as functions of time.
\end{enumerate}

This is primarily orchestrated in \path{diss_lindblad_tail/diss_lindblad/experiment.py}.

\subsection{Liouville-space superoperator methods}
The mathematically most explicit part of the repository is the Liouville-space representation in \path{diss_lindblad_tail/diss_lindblad/lindblad.py}. If $\mathrm{vec}$ denotes column-stacking vectorization, then the code uses
\[
\frac{d}{dt}\lvert \rho \rangle\!\rangle
= \mathcal{L}\,\lvert \rho \rangle\!\rangle,
\]
with
\[
\mathcal{L}
= -i(I\otimes H - H^T\otimes I)
 + \sum_k \gamma_k
\left(
\overline{L_k}\otimes L_k
- \frac{1}{2}I\otimes L_k^\dagger L_k
- \frac{1}{2}(L_k^\dagger L_k)^T\otimes I
\right).
\]
This is a concrete test of superoperator calculus, vectorization identities, and deterministic density-matrix propagation via $\exp(\mathcal{L}t)$.

\subsection{Canonical channels that are compared}
Three dissipative channel families are implemented and contrasted:
\begin{enumerate}[label=\alph*.]
  \item \textbf{Target cooling:}
  \[
  L_k = \lvert \psi_\star\rangle \langle \psi_k^\perp \rvert,
  \]
  where $\{\psi_k^\perp\}$ spans the orthogonal complement of the target.
  \item \textbf{Amplitude damping:}
  local $\sigma^- = \lvert 0\rangle\langle 1\rvert$ operators on each qubit.
  \item \textbf{Dephasing:}
  local $Z$ operators on each qubit.
\end{enumerate}

These are not merely implemented as textbook examples. They serve different mathematical purposes:
\begin{itemize}
  \item target cooling is the \emph{idealized teacher} channel that guarantees convergence to the target;
  \item amplitude damping and dephasing are standard physical baselines not aligned to an arbitrary target;
  \item the comparison asks how much of the target-directed behavior is special to global knowledge of $\psi_\star$.
\end{itemize}

\subsection{Named target states and entanglement structure}
The code explicitly generates:
\begin{itemize}
  \item Haar-random states,
  \item GHZ states,
  \item $W$ states.
\end{itemize}
This matters mathematically because those states probe different entanglement structures. The repository is therefore testing whether dissipative correction behaves similarly for:
\begin{itemize}
  \item generic states,
  \item highly nonlocal cat-like states,
  \item distributed single-excitation states.
\end{itemize}

\subsection{What the tests verify}
The test file \path{diss_lindblad_tail/tests/test_sanity.py} shows what the author regards as the important mathematical invariants:
\begin{itemize}
  \item pure-state density matrices are valid,
  \item unitary channels preserve purity,
  \item zero dissipation leaves the state unchanged,
  \item Lindblad evolution preserves trace,
  \item evolved states remain positive semidefinite,
  \item target cooling converges to the target,
  \item the full untruncated circuit reproduces the target state,
  \item incremental evolution agrees with direct exponentiation.
\end{itemize}
So the package is not only simulating open quantum dynamics; it is also testing numerical faithfulness to the mathematical theory of CPTP semigroups.

\section{Mathematical Ideas in \texttt{experiments/hydrogen\_qalchemy/}}

\subsection{Hydrogenic eigenfunctions as quantum data}
The file \path{experiments/hydrogen_qalchemy/hydrogen_wavefunction.py} introduces a different mathematical domain: the bound-state wavefunctions of the Schr\"odinger hydrogen atom.

The state is modeled in the standard separated form
\[
\psi_{n\ell m}(r,\theta,\phi) = R_{n\ell}(r)Y_{\ell m}(\theta,\phi),
\]
where:
\begin{itemize}
  \item $R_{n\ell}$ is built from generalized Laguerre polynomials and a stable normalization formula using log-gamma functions,
  \item $Y_{\ell m}$ is the spherical harmonic,
  \item the reduced electron-nucleus mass and reduced Bohr radius are included.
\end{itemize}

This means the repository tests more than abstract random states. It tests physically structured wavefunctions with nodal surfaces, angular momentum labels $(n,\ell,m)$, and analytically known radial-angular factorization.

\subsection{Discretization and embedding into a qubit statevector}
The hydrogen wavefunction is sampled on an $x$-$z$ slice with $y=0$, on a square grid whose side length is
\[
2^{n_q/2},
\]
where $n_q$ is the number of qubits and must be even. The sampled complex amplitudes are flattened and normalized into a quantum statevector. This tests the mathematics of:
\begin{itemize}
  \item spatial discretization of a continuous wavefunction,
  \item amplitude encoding of classical scientific data into a quantum state,
  \item compatibility between physical wavefunctions and qubit register dimensions.
\end{itemize}

\subsection{Approximate state preparation through Q-Alchemy}
The script \path{experiments/hydrogen_qalchemy/run_hydrogen_qalchemy.py} sends the target statevector to Q-Alchemy and compiles a preparation circuit. The repository then measures fidelity after each gate in the compiled circuit.

The mathematical idea under test here is:
\begin{quote}
How does an approximate, compiled state-preparation circuit accumulate fidelity to a structured target state as its gate sequence unfolds?
\end{quote}

The code therefore treats the circuit not just as a black box but as a trajectory in state space. It records:
\begin{itemize}
  \item gate-by-gate fidelity values,
  \item final circuit fidelity,
  \item monotonicity ratios for candidate hydrogen targets.
\end{itemize}

This leads to a secondary search problem: some hydrogenic targets may produce more monotone or better-behaved compiled fidelity curves than others.

\subsection{The search over quantum numbers}
The function \path{search_monotonic_targets} loops over admissible $(n,\ell,m)$ and filters states by:
\begin{itemize}
  \item minimum final fidelity,
  \item minimum fraction of non-decreasing gate steps.
\end{itemize}
Mathematically, this is a search over a discrete family of physically meaningful basis states to find states whose compiled preparation trajectories are particularly regular.

\section{Exact Target-Cooling Mathematics}

\subsection{Teacher channel with an analytic fidelity law}
The most elegant mathematical construction in the repository is the exact target-cooling channel in \path{experiments/hydrogen_qalchemy/tail/target_cooling.py}. It uses jump operators
\[
L_k = \sqrt{\gamma}\,\lvert \psi_\star\rangle \langle \psi_k^\perp\rvert,
\qquad k=1,\dots,d-1,
\]
so that the target projector
\[
P = \lvert \psi_\star\rangle \langle \psi_\star\rvert
\]
is the unique steady state.

The code exploits a block decomposition relative to $P$ and $P_\perp = I-P$, which yields an exact closed-form evolution law. In particular, if
\[
F_0 = \langle \psi_\star \rvert \rho_0 \lvert \psi_\star \rangle,
\]
then
\[
F(t) = 1 - (1-F_0)e^{-\gamma t}.
\]

This is important for the whole repository because it provides:
\begin{itemize}
  \item an analytically controlled teacher process,
  \item a monotone target-directed benchmark,
  \item a way to scale to dimensions where a full $d^2\times d^2$ Liouvillian would be infeasible.
\end{itemize}

\subsection{Algorithmic complexity as part of the mathematics}
The code is explicit that a naive Liouvillian is impossible for $d=1024$ and beyond. Instead, the teacher tail is implemented using rank-structured updates and matrix-vector products in $O(d^2)$ time per step. Thus the repository is also testing a mathematical-algorithmic idea:
\begin{quote}
Exploit problem-specific structure of the dissipator instead of treating the open-system evolution as a generic superoperator.
\end{quote}

\subsection{Verified properties}
The target-cooling acceptance tests in \path{experiments/hydrogen_qalchemy/tail/tests.py} confirm the intended mathematics:
\begin{itemize}
  \item the target state is a fixed point,
  \item trace is preserved,
  \item the analytic fidelity law is exact,
  \item fidelity is monotone non-decreasing,
  \item long-time evolution converges to fidelity $1$.
\end{itemize}

\section{Hardware-Constrained Dissipation as an Approximation Problem}

\subsection{From ideal global jumps to realizable local channels}
The file \path{experiments/hydrogen_qalchemy/tail/constraint_dissipation.py} asks the key practical question of the repository:
\begin{quote}
If the ideal target-cooling jump operators are globally nonlocal and unrealistic, what family of physically local or near-local CPTP maps can approximate their effect?
\end{quote}

This is formulated as a model class, \texttt{ConstraintDissipation}, parameterized by rates and optional unitary angles. The mathematical idea is not a single channel, but a \emph{model family} of dissipative dynamics with tunable expressive power.

\subsection{Local channel basis}
The constrained model includes several elementary channels.

\paragraph{Single-qubit channels.}
For each qubit, the repository can include:
\begin{itemize}
  \item amplitude damping,
  \item dephasing,
  \item depolarizing noise.
\end{itemize}

\paragraph{Two-qubit edge channels.}
On each allowed edge of a connectivity graph, it can include:
\begin{itemize}
  \item correlated $ZZ$ dephasing,
  \item Bell-state pumping toward $\lvert \Phi^+\rangle$ or $\lvert \Psi^-\rangle$,
  \item parity pumping in the $ZZ$ or $XX$ basis,
  \item collective decay with jump operator proportional to
  \[
  \sigma^-_i \pm \sigma^-_j.
  \]
\end{itemize}

\paragraph{Ancilla-reset gadgets.}
The code also includes one- and two-qubit ancilla-reset channels obtained by:
\begin{enumerate}[label=\alph*.]
  \item coupling system qubits to a fresh ancilla by a parameterized unitary,
  \item tracing or resetting the ancilla,
  \item interpreting the resulting reduced map as a programmable dissipative primitive.
\end{enumerate}

This is mathematically significant because ancilla-reset gadgets leave the pure Lindblad-generator picture and move into a broader class of discrete CPTP maps.

\subsection{Graph topology as a mathematical constraint}
The class \texttt{ConstraintConfig} allows several connectivity modes:
\begin{itemize}
  \item chain,
  \item custom graph,
  \item all-to-all.
\end{itemize}
Therefore one of the ideas tested in the repository is how geometry limits dissipative control. The difference between chain and all-to-all is not cosmetic; it changes the space of admissible edge channels and therefore the approximation power of the constrained model.

\subsection{Trotterized dissipative composition}
The constrained step is applied as a first-order Trotter composition of exact local CPTP maps. This tests whether a sequence of individually valid local channels can act as a useful approximation to a desired global semigroup. In other words, the repository studies a synthesis problem for open-system dynamics:
\[
\mathcal{E}_{\Delta t}
\approx
\mathcal{E}^{(m)}_{\Delta t}\circ\cdots\circ \mathcal{E}^{(2)}_{\Delta t}\circ \mathcal{E}^{(1)}_{\Delta t}.
\]

\subsection{Kraus-efficient tensor reshaping}
The code applies local channels without constructing full lifted Kraus operators on the entire Hilbert space. Instead, it reshapes $\rho$ into a tensor and contracts only the active indices. This is algorithmically important and mathematically aligned with tensor-network thinking: locality in the channel is translated into low-overhead tensor contractions.

\section{Optimization Problems Tested in the Repository}

\subsection{Teacher-supervised fitting}
The file \path{experiments/hydrogen_qalchemy/tail/fit_constraint_to_target.py} formulates the constrained-tail design problem as supervised approximation to the teacher trajectory.

Two losses are implemented:
\begin{align*}
\mathcal{L}_{\mathrm{Frob}}
&=
\frac{1}{SJ}\sum_{s,j}
\left\lVert
\rho^{\mathrm{model}}_{s}(t_j)
- \rho^{\mathrm{teacher}}_{s}(t_j)
\right\rVert_F^2, \\
\mathcal{L}_{\mathrm{fid}}
&=
\frac{1}{SJ}\sum_{s,j}
\left(
F^{\mathrm{model}}_{s}(t_j)
- F^{\mathrm{teacher}}_{s}(t_j)
\right)^2.
\end{align*}

This means the repository is testing at least two mathematically distinct notions of approximation:
\begin{itemize}
  \item match the full density-matrix trajectory,
  \item match only the target-fidelity curve.
\end{itemize}

\subsection{Direct-to-target optimization}
The file \path{experiments/hydrogen_qalchemy/tail/fit_constraint_direct_to_target.py} removes the teacher and directly optimizes for high target fidelity. Two objective styles appear:
\begin{itemize}
  \item terminal-fidelity loss,
  \item weighted curve loss over the whole dissipative trajectory.
\end{itemize}

So the repository is explicitly comparing two philosophies:
\begin{enumerate}[label=\alph*.]
  \item imitate an analytically ideal channel,
  \item optimize directly for the end task.
\end{enumerate}

\subsection{Parameterization and bounded rates}
Rate parameters are bounded in $[0,\gamma_{\max}]$ using sigmoid or clipping logic, while ancilla-unitary parameters are left unbounded. This is a standard but meaningful mathematical design choice: it separates physically interpretable dissipation strengths from more expressive control angles.

\subsection{Training distributions over states}
The fitting routines do not train on only one input state. They generate mixed training sets including:
\begin{itemize}
  \item the actual circuit output state,
  \item the target state,
  \item random pure states,
  \item computational basis states,
  \item random mixed states.
\end{itemize}
This reveals another tested idea: whether a constrained dissipative model can generalize as an \emph{operator family}, not merely overfit one initial density matrix.

\section{Comparison Experiments and Hypotheses}
The script \path{experiments/hydrogen_qalchemy/compare_tails.py} makes the intended scientific comparison very explicit. It compares:
\begin{itemize}
  \item the teacher tail,
  \item the constrained fitted tail,
  \item optionally a direct-to-target optimized constrained tail.
\end{itemize}

In effect, the repository is testing the following hypotheses.

\subsection*{Hypothesis 1: dissipative post-processing can repair imperfect preparation}
After a compiled state-preparation circuit outputs $\rho_0$, the repository asks whether a dissipative tail can increase
\[
\langle \psi_\star \rvert \rho_0 \lvert \psi_\star \rangle.
\]

\subsection*{Hypothesis 2: the ideal teacher channel has analytically controlled monotone behavior}
The code checks this directly through the formula
\[
F(t)=1-(1-F_0)e^{-\gamma t}.
\]

\subsection*{Hypothesis 3: constrained local channels can approximate the teacher nontrivially}
This is the main approximation problem. The code asks how much of the teacher advantage survives once one restricts to local noise, edge channels, and geometry.

\subsection*{Hypothesis 4: ancilla-reset channels add expressive power}
The presence of \path{experiments/hydrogen_qalchemy/tail/example_edge_ancilla_reset.py} indicates an explicit test that edge ancilla-reset primitives can outperform a do-nothing baseline and possibly close some of the gap to the teacher.

\section{What Role the Q-Alchemy SDK Plays}
The directory \path{q-alchemy-sdk-py/} is mostly not where the new mathematics of this repository lives. Its main role is enabling state preparation through a remote synthesis service. Still, it embodies a few mathematical ideas that matter for interpreting the experiments:
\begin{itemize}
  \item quantum state preparation as an optimization or decomposition problem,
  \item approximate loading with a user-specified fidelity-loss tolerance,
  \item method choices such as hierarchical Tucker, iterative Tucker, swap-pivot, and low-rank variants,
  \item examples based on Schmidt compression and tensor-structured data handling.
\end{itemize}

The most relevant point is that the repository treats state preparation as \emph{approximate and structured}, not exact by default. That assumption is what makes the dissipative-tail experiments meaningful.

\section{Most Important Mathematical Themes, Summarized}
The mathematical ideas tested in this repository can be summarized as follows.

\subsection*{1. Open-system control toward a known pure target}
The repository repeatedly studies target-aware dissipation that attempts to steer an arbitrary or imperfect state toward a specified $\lvert \psi_\star\rangle$.

\subsection*{2. Exact Lindblad theory versus realizable CPTP approximations}
One part of the code uses exact Liouvillian evolution; another replaces it with local Kraus maps and ancilla-reset constructions. The tension between those two descriptions is central.

\subsection*{3. Structured quantum data from physics}
Hydrogenic eigenfunctions give the repository a scientifically structured family of target states rather than only random vectors.

\subsection*{4. Approximate state preparation as a precursor error model}
The code assumes state preparation is imperfect, either because a circuit is truncated or because Q-Alchemy returns an approximate compilation. Dissipation is then studied as a repair mechanism.

\subsection*{5. Optimization over channel families}
The constrained-tail code is effectively learning a dissipative channel family under locality, topology, and bounded-rate constraints.

\subsection*{6. Provable analytic behavior where possible}
The exact target-cooling channel is attractive in this repository precisely because its fidelity law and steady state are analytically controlled.

\section{Conclusion}
The repository tests a coherent mathematical program: use state-preparation circuits to get close to a target state, then use engineered dissipation to move closer. Its most substantive mathematical contributions are the combination of:
\begin{itemize}
  \item Liouville-space Lindblad simulation,
  \item exact target-cooling channels with closed-form fidelity dynamics,
  \item constrained dissipative models built from local and edge CPTP primitives,
  \item optimization-based matching of realizable dissipative tails to ideal target-directed behavior,
  \item physically structured target states derived from hydrogenic wavefunctions.
\end{itemize}

In short, this repository is testing the mathematics of \emph{dissipative correction of approximate quantum state preparation}, with increasing realism from exact global models to topology-limited, hardware-motivated channel families.

\end{document}
